
\documentclass[10pt,letterpaper,unboxed,cm]{article}
\usepackage[margin=1in]{geometry}
\usepackage{graphicx}
\usepackage{enumerate, comment}
\usepackage{adjustbox}
\usepackage{amsmath,mathabx}

\newcommand{\st}{~\mid~}
\newcommand{\ind}{$~~~~~$}




\begin{document}

\hfill{CSE 21 Winter 2026}

\hfill{Homework 2}

\hfill{Due date: Monday, October 12 at 11:59pm}

\begin{center}
\begin{minipage}[t]{5.7in}
\rule{\linewidth}{2pt}
{\sc Instructions}\newline

This homework should be done in {\bf groups} of one to four students, without assistance from anyone besides the instructional staff and your group members.  Homework must be submitted through Gradescope by a {\bf single representative} of your group and received by {\bf 11:59pm} on the due date.\\
 

Students should consult their textbook, class notes, lecture slides, podcasts, group members, instructors, TAs, and tutors when 
they need help with homework. You may ask questions about the homework in office hours, but questions on Piazza should be private, visible only to instructors. 

This assignment will be graded for not only the \emph{correctness} of your answers, but on your ability to present your ideas clearly and logically. You should explain or justify, present clearly how you arrived at your conclusions and justify the correctness of your answers with mathematically sound reasoning (unless explicitly told not to). Whether you use formal proof techniques or write a more informal argument for why something is true, your answers should always be well-supported. Your goal should be to {\bf convince the reader} that your results and methods are sound.






{\sc Key Concepts} Inclusion/Exclusion, Binomial Identities, Stars and Bars.

(Note: For this homework, you can leave your answers in terms of exponentials, factorials, binomial coefficients, etc.)

\rule{\linewidth}{2pt}
\end{minipage} \hfill

\end{center}

\emph{(Note: for justifying counting arguments, a good rule of thumb is to explain how you came up with every term and factor of your answer. You can leave
your answer in terms of exponentials rather than compute the exact numerical value.)}
\begin{enumerate}





\item
A license plate is a string over the alphabet of 36 characters consisting of 26 uppercase latin letters and 10 numerals (0 through 9) (\emph{2 points for correct expression and 2 points for justification})

\begin{enumerate}

\item
(4 points)
How many 7-character license plates start with a letter or end with a numeral?

\textbf{Solution}

We use the Inclusion-Exclusion Principle.

There are
\[ 26(36^6) \]

plates that start with a letter, since there are 26 choices for the first character and 36 choices for the rest.

There are
\[ 10(36^6) \]

plates that end with a numeral, since there are 10 choices for the last
character and 36 choices for the rest 6 characters. 

Plates that start with a letter and end with a numeral are counted twice.
There are
\[ 26(36^5)(10) \]
such plates.

Therefore, by the Inclusion-Exclusion Principle, the answer is
\[ \boxed{26(36^6) + 10(36^6) - 26(10)(36^5)} \]

\item
(4 points)
How many 7-character license plates use at least one A, at least one B and at least one numeral?

\textbf{Solution}

We can use the Inclusion-Exclusion Principle and De Morgan's Law her.

Let $A$ be the set of plates with at least one A, $B$ be the set of plates with at least one B, and $C$ be the set of plates with at least one numeral.

We want to find
\[ |A \cap B \cap C| \]

By De Morgan's Law,
\[ |A \cap B \cap C| 
=
36^7 - |A' \cup B' \cup C'|\]

Using the Inclusion-Exclusion Principle,

\[ |A' \cup B' \cup C'|
= |A'|+|B'|+|C'| -|A'\cap B'| -|A'\cap C'| -|B'\cap C'| +|A'\cap B'\cap C'|\]

The individual terms ar
\[ |A'|=35^7,\qquad |B'|=35^7,\qquad |C'|=26^7, \]

since excluding A or B leaves 35 choices for each position, while
excluding all numerals leaves 26 choices for each position.

For the intersections,
\[ |A'\cap B'|=34^7, \]
since excluding both A and B leaves 34 choices for each position.

Also,
\[ |A'\cap C'|=25^7,\qquad |B'\cap C'|=25^7, \]
since the plate can only contain letters, with either A or B excluded.

Finally,
\[ |A'\cap B'\cap C'|=24^7,\]
since the plate can only contain letters, with both A and B excluded.

Therefore, the answer is
\[ \boxed{ 36^7-\left(2(35^7)+26^7-34^7-2(25^7)+24^7\right) }
\]

\item
(4 points)
How many 7-character license plates contain only numerals that all add up to 10?

\textbf{Solution}

Let $x_1,x_2,\dots,x_7$ be the seven numerals. We need
\[x_1+x_2+\dots+x_7=10,\]
where each $x_i$ is between 0 and 9.

For the base case, by Stars and Bars the number of
nonnegative integer solutions is

\[\binom{10+7-1}{7-1}=\binom{16}{6}\]

However, this includes the 7 cases where one numeral is 10 and all
the others are 0. These are not valid numerals.

Therefore, the answer is
\[ \boxed{\binom{16}{6}-7}\]


\item
(4 points)
How many 7-character license plates are there that are all letters and they are in alphabetical order (repeats allowed)?

\textbf{Solution}

Since the letters must be in alphabetical order, once we choose how
many times each letter appears, their order is fixed.

Let $x_1,x_2,\ldots,x_{26}$ be the number of times each letter appears.

Then
\[x_a+x_b+\dots+x_{z}=7,\]
where each $x_i \geq 0$ since repeats are allowed.

By Stars and Bars, the number of solutions is
\[\binom{7+26-1}{26-1}=\boxed{\binom{32}{25}}\]

\end{enumerate}

\item
An art museum with 12 different rooms just got a shipment of 9 new (identical) prints (\emph{2 points for correct expression and 2 points for justification})
\begin{enumerate}
\item
(4 points)
How many ways can the museum curator arrange all 9 prints among the 12 rooms in the museum?

\textbf{Solution}

Let $x_i$ be the number of prints placed in room $i$. Then
\[x_1+x_2+\dots+x_{12}=9,\]
where each $x_i \geq 0$.

There are 9 identical prints and 12 different rooms. By Stars and Bars,
the number of ways to distribute the prints is
\[\binom{9+12-1}{12-1}=\boxed{\binom{20}{11}}.\]

\item
(4 points)
How many ways can the museum curator arrange all 9 prints among the 12 rooms in the museum so that each print is in a different room?

\textbf{Solution}
Since each print must be in a different room, exactly 9 of the 12 rooms
will contain one print.

We choose which 9 rooms will contain a print:
\[\binom{12}{9}\]

Since the prints are identical, there are no additional ways to arrange
them once the rooms are chosen.

Therefore, the answer is
\[\boxed{\binom{12}{9}}\]

\item
(4 points)
How many ways can the museum curator arrange all 9 prints among the 12 rooms in the museum so that at least 3 of the rooms are not empty?

\textbf{Solution}

We use the Complement Principle here.

The total number of ways to arrange the 9 prints among the 12 rooms is
\[\binom{20}{11}\]

The complement of having at least 3 nonempty rooms is having exactly
1 or exactly 2 nonempty rooms.

For exactly 1 nonempty room, we choose 1 of the 12 rooms:
\[\binom{12}{1}.\]

For exactly 2 nonempty rooms, we first choose the 2 rooms:
\[\binom{12}{2}\]

The 9 prints can be distributed between these 2 rooms in
\[\binom{10}{1}-2\]
ways. We subtract 2 because the cases where either room gets all 9
prints would be invalid, since we are counting for non-empty rooms.

Therefore, the answer is
\[\boxed{\binom{20}{11}-[\binom{12}{1}+\binom{12}{2}\left(\binom{10}{1}-2\right)]}\]

\item
(4 points)
How many ways can the museum curator arrange all 9 prints among the 12 rooms such that at least one room has at least 3 prints?

\textbf{Solution}

The total number of ways to distribute the 9 prints among the 12 rooms is
\[\binom{20}{11}\]

The complement is that every room has at most 2 prints. The possible
distributions of the 9 prints are
\[1+1+1+1+1+1+1+1+1,\]

\[2+1+1+1+1+1+1+1,\]

\[2+2+1+1+1+1+1,\]

\[2+2+2+1+1+1,\]
and
\[ 2+2+2+2+1\]

For these five cases, we choose which rooms contain 2 prints and which
rooms contain 1 print. This gives
\[
\binom{12}{9}
+\binom{12}{1}\binom{11}{7}
+\binom{12}{2}\binom{10}{5}
+\binom{12}{3}\binom{9}{3}
+\binom{12}{4}\binom{8}{1}
\]

Therefore, by the Complement Principle, the answer is
\[
\boxed{\binom{20}{11} -
[\binom{12}{9}
+\binom{12}{1}\binom{11}{7}
+\binom{12}{2}\binom{10}{5}
+\binom{12}{3}\binom{9}{3}
+\binom{12}{4}\binom{8}{1}
]}\]


\item
(4 points)
How many ways can the museum curator arrange all 9 prints in 3 of the 12 rooms such that each room has exactly 3 prints? 

\textbf{Solution}

Since all 9 prints must be placed in exactly 3 rooms, with each room
containing exactly 3 prints, we only need to choose which 3 rooms are used.

There are
\[\binom{12}{3}\]
ways to choose these rooms.

There are no additional ways to arrange them because the prints are identical.

Hence, the answer is

\[\boxed{\binom{12}{3}}\]

\end{enumerate}
  






\item (12 points)
Compute the number of \emph{integer solutions} for each equation:

(Include a justification)

\begin{enumerate}
\item
$$a_1+a_2+a_3+a_4 + a_5 + a_6 = 30$$
$1\leq a_1$\\
$2\leq a_2$\\
$3\leq a_3$\\
$4\leq a_4$\\
$5\leq a_5$\\
$6\leq a_6$

\textbf{Solution}

We first assign the minimum required value to each variable. The minimum total is
\[1+2+3+4+5+6=21\]

This leaves
\[30-21=9\]
to distribute among the 6 variables.

Let
\[x_1=a_1-1,\quad x_2=a_2-2,\quad x_3=a_3-3,\]
\[x_4=a_4-4,\quad x_5=a_5-5,\quad x_6=a_6-6\]

Then each $x_i \geq 0$ and
\[x_1+x_2+x_3+x_4+x_5+x_6=9\]

By Stars and Bars, the number of nonnegative integer solutions is
\[\binom{9+6-1}{6-1}=\boxed{\binom{14}{5}}\]

\item
$$a_1+a_2+a_3+a_4 + a_5 + a_6= 30$$
At least one of the variables is less than or equal to 3.

\textbf{Solution}\\
The total number of nonnegative integer solutions to
\[a_1+a_2+a_3+a_4+a_5+a_6=30\]
is, by Stars and Bars,
\[\binom{30+6-1}{6-1}=\binom{35}{5}\]

The complement is that every variable is at least 4. Assigning 4 to
each of the 6 variables uses
\[6(4)=24,\]
leaving
\[30-24=6\]

By Stars and Bars, the number of solutions in the complement is
\[\binom{6+6-1}{6-1}=\binom{11}{5}\]

Therefore, the number of solutions where at least one variable is less
than or equal to 3 is
\[\boxed{\binom{35}{5}-\binom{11}{5}}\]

\item
$$a_1+a_2+a_3+a_4 + a_5 + a_6 = 30$$
$0\leq a_1\leq 7$\\
$0\leq a_2\leq 7$\\
$0\leq a_3\leq 7$\\
$0\leq a_4\leq 7$\\
$0\leq a_5\leq 7$\\
$0\leq a_6\leq 7$

\textbf{Solution}

We want to count the solutions to
\[a_1+a_2+a_3+a_4+a_5+a_6=30, \
0\leq a_i\leq 7.\]

Ignoring the upper bounds, by stars and bars there are
\[\binom{30+6-1}{6-1}=\binom{35}{5}\]solutions.

Now use inclusion-exclusion to remove solutions where some
$a_i\geq 8$.

If one variable is at least $8$, choose that variable in
\[\binom{6}{1}\]
ways. Subtract $8$ from that particular variable. The remaining amount to
distribute among the rest is
\[30-8=22\]

By stars and bars, the number of solutions we get
\[\binom{22+6-1}{6-1}=\binom{27}{5}\]

Thus, we subtract
\[\binom{6}{1}\binom{27}{5}\]

If two variables are at least $8$, choose those two variables in
\[\binom{6}{2}\]
ways. Subtracting $8$ from both of them makes the remaining sum
\[30-16=14\]

By stars and bars, the number of solutions is
\[\binom{14+6-1}{6-1}=\binom{19}{5}\]

Thus, we add back
\[\binom{6}{2}\binom{19}{5}\]

If three variables are at least $8$, choose those three variables in
\[\binom{6}{3}\]
ways. Subtract $8$ from each of them. The remaining amount is
\[30-24=6\]

By stars and bars, the number of solutions is
\[\binom{6+6-1}{6-1}=\binom{11}{5}\]

Thus, we subtract
\[\binom{6}{3}\binom{11}{5}\]

Four variables cannot each be at least $8$ since
\[4(8)>30\]

Therefore, by inclusion-exclusion, the number of valid solutions is
\[\boxed{\binom{35}{5}
-\binom{6}{1}\binom{27}{5}
+\binom{6}{2}\binom{19}{5}
-\binom{6}{3}\binom{11}{5}}\]

\end{enumerate}
\item
\begin{enumerate}

\item (8 points)
For $n\geq 0$ and $k\geq 0$ consider the identity:

$$k{n\choose k} = n{n-1\choose k-1}$$


Prove it combinatorially by counting the same set in two different ways. 

\textbf{Solution}

Consider choosing $k$ balls from $n$ distinct balls and marking one of the chosen balls.

One way is to first choose the $k$ balls. There are
\[\binom{n}{k}\]

ways to do this. Then, there are $k$ choices for which of the chosen balls to mark. Therefore, the number of ways is

\[k\binom{n}{k}\]

Another way is to first choose the marked ball. There are $n$ choices.
Then, we choose the remaining $k-1$ balls from the remaining $n-1$ balls.

There are
\[\binom{n-1}{k-1}\]
ways to do this. Therefore, the number of ways is
\[n\binom{n-1}{k-1}\]

Since both methods count the same set of $k$ chosen balls with one marked ball, we have
\[\boxed{k\binom{n}{k}=n\binom{n-1}{k-1}}\]


\item (8 points)
For $n\geq 3$, prove that:

$${n+1\choose 3} = 1(n-1) + 2(n-2) + 3(n-3) + \dots + (n-3)3 + (n-2)2 + (n-1)1
$$

using a combinatorial argument.

[Hint: count the fixed density binary strings with 3 ones based on where is the middle one)]

\textbf{Solution}

Consider binary strings of length $n+1$ with exactly three $1$'s.
We can choose the three positions containing the $1$'s in

\[\binom{n+1}{3}\]
ways.

Now count the same strings based on the position of the middle $1$.
If the middle $1$ is in position $2$, there is $1$ choice for the
first $1$ and $n-1$ choices for the last $1$, giving

\[1(n-1)\]

If the middle $1$ is in position $3$, there are $2$ choices for the
first $1$ and $n-2$ choices for the last $1$, giving
\[2(n-2)\]

Continuing this way, placing the middle one is all possible position, we get
\[1(n-1)+2(n-2)+3(n-3)+\dots+(n-1)(1)\]
Since both expressions count the same set of binary strings,
\[\boxed{\binom{n+1}{3}=1(n-1)+2(n-2)+3(n-3)+\dots+(n-1)(1)}\]

\end{enumerate}

\end{enumerate}

\end{document}